\begin{abstract}
Estimating grape crop before harvest still relies largely on sparse, destructive hand sampling, and cameras only help if they can see the fruit. In vineyards, leaves hide much of it from any single view. Robotic harvesters handle occlusion with active perception, moving a camera to a better viewpoint for each fruit. A vehicle driving down a straight vine row cannot do this, but it already supplies a dense sweep of lateral viewpoints as a by-product of its motion. We ask whether that passive sweep, captured by one ordinary RGB camera, can recover hidden fruit well enough to improve a downstream task: counting grapes. Our pipeline segments fruit and foliage with SAM~3, registers views using SIFT features found only on fruit, solves one similarity transform per view jointly with held-out validation, and composites only fruit colours that agree across several views spanning a minimum baseline. It never synthesises unseen fruit. Using a phone on a laboratory rail as a surrogate for a vehicle, we vary viewpoint spacing (20--60~mm) and occlusion (3--6 foreground branches). We evaluate in two ways: pixel-wise against a leaf-free photograph, and by count, using GrapeSAM, a point model trained on the ViViD-5K vineyard dataset and left untuned on our data, against 25 hand-labelled grapes. At 20~mm spacing with five branches, fruit pixel recall rose from 42\% to 87\% and count F1 from 0.42 to 0.77. Benefits fell as sampling became sparser and occlusion heavier. At six branches, recovery was partial (25\% to 42\% recall), and at 60~mm it fell below the single view. All conditions share one physical scene, so this is a controlled proof of concept, not an estimate of field accuracy or yield.
\end{abstract}